\section*{ЗАКЛЮЧЕНИЕ}
\addcontentsline{toc}{section}{ЗАКЛЮЧЕНИЕ}

В результате выполнения научно-исследовательской работы по теме
``Методика обнаружения несанкционированных действий и аномальной
активности в корпоративной сети с использованием нейронных сетей''
получены следующие основные результаты.

\begin{enumerate}
    \item \textbf{Теоретический анализ.} Проведён систематический
    обзор методов и технологий обнаружения аномалий и атак: классические
    статистические подходы, методы машинного обучения, восемь типов
    нейросетевых архитектур. Систематизированы методы моделирования
    атак (шаблонные, генеративные, adversarial-RL). Выполнен
    сравнительный анализ десяти публично доступных датасетов NIDS по
    восьми формализованным критериям; в качестве основного выбран
    UNSW-NB15, вспомогательного --- CICIDS2017. Все используемые
    числовые результаты привязаны к 45 верифицированным первоисточникам,
    включая публикации IEEE Xplore, ACM, NDSS, NeurIPS, ICML, RFC и
    действующие государственные стандарты.

    \item \textbf{Проектирование методики.} Сформулирована методика
    обнаружения аномальной активности, включающая: подсистему
    предобработки flow-данных (очистка, кодирование, log-преобразование,
    robust-scaling, скользящие окна $W = 32$, $S = 8$, last-агрегация
    меток); нейросетевой детектор на архитектуре CNN-LSTM;
    подсистему активного тестирования --- ИИ-злоумышленник в форме
    \emph{templates + FSM + drift-injector}; блок drift-мониторинга
    на основе индексов PSI, KL и MMD; pipeline формирования алертов с
    пятиступенчатой шкалой severity и временной дедупликацией.

    \item \textbf{Программная реализация.} Реализован полный
    программный прототип \texttt{diploma\_nids} объёмом порядка
    4230 строк Python (PyTorch, scikit-learn, XGBoost, FastAPI,
    Streamlit) с восемью подпакетами, 26 конфигурационными
    YAML-файлами, 12 скриптами и 9 модулями юнит-тестов.
    Реализованы 15 моделей-кандидатов в едином registry, библиотека
    из 7 параметризованных шаблонов атак, конечный автомат на 11
    состояниях, drift-инжектор с тремя типами дрейфа, REST-сервис
    инференса и Streamlit-дэшборд мониторинга.

    \item \textbf{Экспериментальное исследование.} Проведены шесть
    экспериментов (E1--E6) с 22 запусками и 14 моделями в
    идентичных условиях. Установлено: на табличных flow-данных
    UNSW-NB15 лидером по F1 является XGBoost ($0{,}919$), Random
    Forest --- $0{,}917$, ведущие DL-модели (TCN, BiLSTM,
    LSTM/GRU) --- в диапазоне $0{,}82$--$0{,}84$. CNN-LSTM
    показал среднее $F_1 = 0{,}770$, $\text{PR-AUC} = 0{,}955$
    при латентности $0{,}96$\,мс/окно (CPU). Калибровка под
    целевой $\text{FPR} = 0{,}01$ обеспечивает
    $\text{Precision} \approx 0{,}99$; drift-monitor корректно
    срабатывает на всех контролируемых режимах дрейфа
    ($\text{PSI} > 0{,}25$).

    \item \textbf{Подтверждение гипотез.}
    \begin{itemize}
        \item Основная гипотеза $H_0^{\text{main}}$ о нейросетевом
        подходе \emph{уточнена}: на табличных flow-данных классические
        ансамбли (XGBoost, RF) остаются сильным baseline, однако
        CNN-LSTM сохраняет конкурентоспособность по PR-AUC при
        существенно лучшей латентности (0{,}96 мс против 51 мс у RF)
        и единственная среди ведущих моделей моделирует
        последовательную структуру flow-окон. Это согласуется с
        обзорами Ferrag~2020 и Ahmad~2021.
        \item Дополнительная гипотеза $H_0^{\text{aux}}$ о роли
        ИИ-злоумышленника \emph{подтверждена}: ground-truth по 11
        состояниям FSM позволила измерить per-state recall и FPR
        (Эксперимент E4) --- аналитика, недоступная в статических
        датасетах.
        \item Гипотеза переносимости $H_0^{\text{transfer}}$
        \emph{подготовлена} к проверке: датасет CICIDS2017 распакован,
        mapping общих признаков зафиксирован.
    \end{itemize}

    \item \textbf{Соответствие требованиям ВКР.} Все девять задач
    исходной постановки реализованы: изучение методов NIDS/NTA,
    моделирование атак с ИИ, формулировка гипотезы и требований,
    выбор архитектуры нейронной сети, сбор и подготовка данных,
    обучение и сравнение моделей, разработка ИИ-злоумышленника и
    тестирование в связке агент$\to$детектор, реализация прототипа,
    оценка эффективности и доработка. Все четыре формы ожидаемых
    результатов (теоретический анализ, разработанная методика,
    экспериментальные результаты, итоговые материалы) представлены
    в составе настоящего отчёта и сопутствующих артефактов.
\end{enumerate}

\textbf{Цель работы достигнута.} Разработана методика обнаружения
несанкционированных действий и аномальной активности в корпоративной
сети на основе нейронных сетей с активным тестированием посредством
ИИ-злоумышленника. Методика подкреплена воспроизводимым программным
прототипом и экспериментально проверена на реальном датасете
UNSW-NB15.

\textbf{Направления дальнейшего развития:}
\begin{itemize}
    \item Калибровка шаблонов ИИ-злоумышленника под маргинальные
    распределения UNSW (KS-тест и MMD на режиме NORMAL) для
    устранения out-of-distribution shift, зафиксированного в E5.
    \item Полный прогон cross-evaluation на CICIDS2017 с дообучением
    на 10\,\% выборки.
    \item Расширение методики на pcap-уровень с использованием полного
    UNSW-NB15\_1..4.csv ($\sim 2{,}5\,\text{млн}$ записей) для
    использования естественной временной структуры данных.
    \item Реализация cGAN-составляющей агента (WGAN-GP с условием по
    \emph{attack\_cat}) для повышения вариативности синтетических
    сценариев.
    \item Поддержка дополнительных датасетов (TON\_IoT, Bot-IoT,
    LITNET-2020) для проверки переносимости методики на смежные
    сегменты сетей.
    \item Online-обучение с использованием механизмов EWC или
    replay-buffer для адаптации к concept drift без полного
    переобучения.
\end{itemize}
